% !TeX root = ../papers/nc-p3.tex
\section{環状四次曲線を保つ反転窓の特殊化}
\label{sec:polygon-apolar-limits}
\begin{lemma}[特殊ファイバーの積と極限の六式]\label{lem:mixed-rank-loss}
楕円鉛筆の付値環上の特殊化で、一般ファイバーの二束の積の切断空間は八次元とする。
第一・第二座標格子を$A,B$、隣接の極限六式に双対な商写像を$\beta$とする。
両射影写像を解消した曲線族の特殊ファイバーの被約成分上へ積を制限した写像を
$\Phi:A_0\otimes B_0\to\bigoplus_i H^0((LM)|_{P_i})$とし、その階数を$r$とする。
$\Phi$で零になる任意の部分空間の$\beta$による像の次元は$8-r$以下である。
$r=8$なら$\beta$は$\Phi$の像を二平面で割る写像である。
\end{lemma}
\begin{proof}
分数体上で全積が張る八次元空間の中に$A\otimes B$の像の格子$D'$を取り、
さらに一般鉛筆で割った空間での像の格子$H'$を取る。ともに自由で階数は8,6である。
$A\otimes B\to D'$の核の元は一般ファイバーで零なので、両写像を延ばした整な全空間で
切断として零であり、特殊ファイバーの被約成分上でも零である。
従って$D'_0$から$\operatorname{im}\Phi$への全射がある。
$D'_0\to H'_0$も全射で、六式の一定階数の飽和極限に対応する。結論が従う。
\end{proof}

\begin{lemma}[二つの被約モデルの成分]\label{lem:nodal-core-resolution}
一般点で同定された二つの四次完全交叉曲線族について、一方の特殊ファイバーを節点だけを持つ
有理曲線の多角形（整な節点曲線も含む）、他方も被約とする。
両射影を持つ正則な共通解消の特殊ファイバーの被約支持は、有理曲線の一つの環とそれに付く木からなる。
第一射影の正次数成分は環上にある。他方も多角形ならその正次数成分も全てこの環上にあり、
楕円多重点を唯一の特異点とする有理枝の合併なら環全体が一点に写る。
\end{lemma}
\begin{proof}
節点の局所平滑化の完備方程式は$uv=t^a$である。実際、特殊ファイバーの二次項が非退化なので
平行移動と次数ごとの座標変換で$uv$と底の関数にでき、底の単元を吸収できる。
解消で二枝の間に有理二次曲線の鎖を挿入する（$u,t$のイデアルのブローアップの一方のチャートは
$u'v=t^{a-1}$、他方は正則）。必要なら先に有限基底変換し、自己節点も二成分間の交点にできる。
残る第二射影の有理写像は正則曲面上の点のブローアップで解消する。
交点でのブローアップは辺の細分、平滑点では末端の付加である。
各極限曲線族は完全交叉でCohen--Macaulay、一般ファイバーと特殊ファイバーの一般点で正則なので
全空間は正規である。共通解消からの固有双有理射は連結ファイバーを持つ。
特殊ファイバーの因子の押し出しは値域の特殊ファイバー因子である
（主因子$\operatorname{div}(t)$の押し出し）。被約性により、各像成分を支配する成分はただ一つで、
その重複度と写像次数はいずれも1である。対応する正規化間の写像は同型である。

第二像が多角形なら、各節点の二つの原像は連結ファイバーにより、直接の交点または
定値成分の鎖でつながる。各正次数成分の二端を順にたどると単一の環を作る。
異なる節点のファイバーは交わらず、整な自己節点の場合も一つの鎖が両端を結ぶので同じである。
第二像が多重点型で環上に正次数成分が一つだけあれば、環の残りを通じてその異なる二点が同じ像を持ち、
その枝の埋込み性に反する。二つ以上なら相隣る像は特異点で交わるため、各正次数成分の
環上の二端がその枝上の同一の原点に写り、やはり矛盾する。
整な尖点の場合も正規化写像の閉点での単射性で同じである。
\end{proof}

\begin{theorem}[二つの環状モデルの次数一致]\label{thm:two-polygon-specialization}
楕円鉛筆の特殊化で両側の極限二次式が被約な節点多角形を切るとする。
特殊窓が標準ASなら、二つの曲線は同じ多角形として同定でき、
その上の等しい正の多重次数の二束と基点自由鉛筆の表示を持つ。反対階数は6または8となる。
\end{theorem}
\begin{proof}
前補題の環に制限する。共に定値の木は特殊切断の値を制限しない。
それぞれの正次数成分の次数と束の貼り合わせを引き戻した環上の束を$L,M$、多重次数を$d,e$とする。
全切断は各四次元で、元のそれぞれの四座標がその全体を張る。
実際、多角形の零次数鎖の定数を消去すると元の束の切断に戻る。指定した座標は
被約完全交叉を埋め込む全一次切断であり、成分上への引き戻しは単射である。
共通の零次数成分を除くと$LM$の切断は八次元である。
$d=e$なら既出の正次数での積の全射性により補題\ref{lem:mixed-rank-loss}で$r=8$となる。
従って普通の二次元鉛筆の表示となる。基点があれば純テンソル関係を生じて標準ASと両立しない。
二曲線は同じ鎖を縮めた多角形であり、正則性・反対階数・復元は多角形の鉛筆定理による。

$d\ne e$とする。環上の貼り合わせ係数を共に1に正規化できる。
各成分の両端を固定する座標の倍率は二つの貼り合わせスカラーに指標$d,e$で作用するからである
（単項式の二端での係数を取る）。$d,e$は総次数が同じで異なるので有理数上独立で、
代数閉体上その二指標の写像は二次元トーラスへ全射である。
従って切断の積の階数下界は各多重次数ごとに整数行列で検査できる。

詳細な有限検査を示す。各成分で斉次単項式の係数を取り、各節点の両端係数を同じ変数にする。
零次数成分の両節点の変数も同じにする。内部係数は個別の変数で、
各束につき独立な四変数となる。各成分で二多項式を掛け、各係数を16個の双線型単項式で表す。
得た整数行列の二元体での階数$h$は$r$の下界である。
相補鎖への制限で積が零になる二切断空間の次元を$u,v$とする。
$\min(u,v)\ge1$かつ$\max(u,v)\ge 8-h+2$なら、補題\ref{lem:mixed-rank-loss}で
小さい側の一文字を固定すると大きい側にはなお二次元以上の評価核が残る。
転置の次窓と合わせると関係が消える一次元・二次元の交互商で指数成長する。

\path{computations/verify_polygon_multiplication.py}は四単位ずつの全配分を列挙し、
共通巡回移動・反転・二束交換で重複を除く。異なる配分174通りのうち、$h=7$の24通りは全て
上の次元不等式を満たす。$h=8$の150通りでは148通りが満たし、残る二型は
命題\ref{prop:semipositive-pencil-balance}の二変数の商で排除した二型である。
この二型では実際の六式も$\Phi$の像を経由するのでその排除を適用できる。
従って$d\ne e$は標準ASにならない。
\end{proof}

\begin{proposition}[多角形と楕円多重点を組み合わせた二側の排除]\label{prop:polygon-additive-exclusion}
楕円鉛筆の特殊化で、片側の極限二次式が有理多角形を、他側の極限二次式が
被約な楕円多重点型の四次曲線を定めれば、特殊窓は標準ASでない。
\end{proposition}
\begin{proof}
補題\ref{lem:nodal-core-resolution}により多角形の正次数は環に、多重点曲線の正次数は木に載る。
環は第二像の特異点に写る場合をまず扱う。第二像の各枝がどの木に属するかに応じ、
その木の付け根における第一像の点を$x_i$とする。第二枝では第一束の切断は$x_i$での値に等しい。
各線束の局所枠の定数倍は以下の次元に影響しない。
第二束の特異点で消える切断の空間$B'$は三次元である。環上は第二束の値が同一なので、
全積の環への制限は四次元である。$A_0 B'$の積は環で零であり、各第二枝では
$b_i$に$a(x_i)$を掛けた値である。
$m$枝の二次以上の係数は$4-m$個独立であり、一次係数は全係数非零の一式の核である。
$a$が全ての$x_i$で非零となるよう取れば、この三次元は積の像に入るので$r\ge7$である。
$x_i$が全て同一なら、そこを消す第一束の三次元空間と$B'$の積は全成分で零である
（両射影が定値の鎖でも零になることは各端で評価すればよい）。
実際の積の補正は一次元以下だから、一文字を固定して二次元の評価核が残る。

異なる$x_i$があれば、二つの第一切断の値の比が全枝で同じでないようにできる。
一次係数の超平面に非零の値を対角に掛けた二つの像は異なるので、合わせて全$m$座標を張る。
従って$r=8$である。次数2以上の第二枝があれば、その枝だけに支えを持つ切断と、
$x_i$で消える第一切断三次元の積が零である。全枝が次数1なら、二枝のみに支えを持つ
非零切断と、二つの付け根で消える第一切断二次元以上の積が零である。
どの場合も実際の六式で矩形の零点を生じ、指数成長する。

環の第二像が滑らかな点$y$に落ちる場合も含める。
特異点への各枝を含む木が二本以上なら、それらの特異点の原像を結ぶ道は環を通るため、
その点のファイバーの連結性に反する。ただ一つの木に全ての第二正次数成分があれば、
第一束の値はそれら全てで同一の$x$での評価となる。
$y$で消える第二束の三次元空間と$x$で消える第一束の三次元空間の積は零である。
実際、正次数のない他の木は$y$に写り、唯一の木は$x$に写る。
環上の切断とこの木の切断から再び$r\ge4+3$が従い、補正があっても排除される。
\end{proof}

\begin{corollary}\label{cor:reduced-polygon-window-test}
隣接階数6、反対階数8の両対称テンソルで、両反対核が被約四次完全交叉を切り、
一方は整な節点曲線または可約な節点多角形とする。反転窓が標準ASとなる必要十分条件は、
関係六平面が純テンソルを含まず、両曲線の同型のうち二つの埋込み束の多重次数を一致させるものの
グラフへの制限をテンソル汎関数が経由することである。そのとき順序付き鉛筆データを復元できる。
\end{corollary}
\begin{proof}
必要性は両モデルの被約曲線型を分け、上述の二側の定理と、他側が滑らかな場合には
系\ref{cor:one-smooth-apolar-as}を用いればよい。十分性ではグラフ上で$N=LM$として、
等しい正の多重次数だから$H^0(L)H^0(M)=H^0(N)$である。
与えられた汎関数は$H^0(N)$上の対称対合を誘導し、階数6なので核は二次元である。
この核を$W$とすれば第一窓は$H^0(N)/W$への積で定まり、基点は純テンソル関係を与えるため存在しない。
整曲線または多角形の全次数定理と復元定理が適用できる。
\end{proof}
